\documentclass{zhvt-classic}

% 使用本地源流明体（GenRyuMin2 TW）按文件路径加载，不依赖系统字体名
% 版框（内框／界栏／书口侧栏／鱼尾）改用与版心印章同一印泥棕色：
%   印章图片 yshs.png 实色 = RGB(135,68,52) = HTML 874434
\usepackage{xcolor}
\usepackage{tikz}
% sealbrown 双注册：xcolor 供 TikZ/pgf，L3 供 cls 版框
\definecolor{sealbrown}{HTML}{874434}
\ExplSyntaxOn
\color_set:nnn { sealbrown } { HTML } { 874434 }
\ExplSyntaxOff

\zhvtset{
  grid color       = sealbrown,
  main font        = GenRyuMin2TW-SB,
  main font path   = C:/Users/Administrator/texmf/fonts/opentype/public/GenRyuMin2TW/,
jiazhu font      = GenRyuMin2TW-M,
  jiazhu font path = C:/Users/Administrator/texmf/fonts/opentype/public/GenRyuMin2TW/,
  % 页码框容纳 4 字：正常页「九十七」+前置空 1 格=4 行恰好；
  % 100+ 页「一百零一」共 5 行超出，末字「一/二/三」被框裁掉（2026-10-04 实测），故放宽到 5
  shukou page number chars = 5,
  % 书口书名框高从默认 6 字减到 4 字，把省出的 2 字高拨给 chapter 框，
  % 使 6 字书名+卷名放进 9 字高的 chapter 框后不再压线（2026-10-05）
  shukou title chars = 4,
}

\title{史记}[南越列传]
\author{汉·司马迁}
\maker{闲叟漫记·重製}


% ===== 索隐述赞样式（2026-11-08 四次定稿 · 愚大人亲手改定参数）=====
% ⚠ 竖排前提：自定义字体家族必须带 RawFeature={vertical:+vert}；
%   simhei.ttf 加 +vhal 会整体加载失败静默回退，只写 +vert；simfang 同理。
\newCJKfontfamily{\shuzanfont}[Path=C:/Windows/Fonts/,RawFeature={vertical:+vert}]{simhei.ttf}
\newCJKfontfamily{\shuzanbody}[Path=C:/Windows/Fonts/,RawFeature={vertical:+vert}]{simfang.ttf}
% 题头：152x30pt 圆角 6pt 棕底白字；baseline (0,2) 使整框落入界行格子不压版框线；
% 白字节点 x 不补偿、y +8.2pt（像素级墨迹实测 dx/dy≈0，L/R 留白约 20pt）。
\newcommand{\shuzanhead}[1]{%
  \begin{tikzpicture}[baseline={(0,2)}]
    \fill[sealbrown, rounded corners=6pt]
      (-76pt,-15pt) rectangle (76pt,15pt);
    \node[text=white, inner sep=0pt, anchor=center] at (0pt,8.2pt)
      {\shuzanfont\fontsize{23}{23}\selectfont #1};
  \end{tikzpicture}%
}
% ★ 修正书口卷名位置：chapter 框起始位置改回 cls 默认值，避免全局影响。
%   只在每个卷首页通过 \juan 命令里的空命令微调该页位置（安全，不会跨页污染）。
\ExplSyntaxOn
\AtBeginDocument{
  \dim_set:Nn \l__zhvt_shukou_chapter_title_start_dim
    { \l__zhvt_shukou_title_height_dim + \l__zhvt_font_size_dim * 2 }
  \dim_set:Nn \l__zhvt_shukou_page_number_start_dim
    { \l__zhvt_shukou_chapter_title_start_dim + \l__zhvt_shukou_chapter_title_height_dim }
}
\ExplSyntaxOff

\begin{document}

% ★★ 版心（书口）排布覆写：把书名放到鱼尾之下（2026-10-04）
%
%   cls 默认版心自上而下：书名（title 静态框，在鱼尾之上）→ 鱼尾 → 篇名（chapter 动态框）→ 页码
%   本模板改为：          鱼尾（位置不动）→ 书名 → 篇名 → 页码分割线 → 页码
%
%   手法：把 cls 放在 title 静态框里的书名，改放进 chapter 动态框（鱼尾之下）；
%         title 静态框留空。书名取 \l__zhvt_title_str，即 \title{...} 的第一参数，
%         所以换书时只改 \title，不必在此重复写书名。
%
%   ⚠ 三条必须遵守的写法（2026-10-04 实测）：
%     ① 汉字必须用 \hbox{...} 包住。flowfram 会把框内容装进一个宽 14.4mm（= 一个 line_sep
%        = 书口栏宽）的纵向 minipage 并按该宽度折行，而整页靠 PDF /Rotate 90 显示、
%        框内文字本身不变向 —— 不包 \hbox 时「围炉夜话」4 字会被逐字折成 4 行，
%        在版面上就成了横排 4 字并横向溢出版框。
%     ② 竖排下**一个 \hbox 就是一整列**。书名与篇名要同处一列，必须放进同一个 \hbox；
%        分成两个 \hbox 时，后一个会被当成新的一列、甩到隔壁栏去
%        （实测「序」跑到书口栏外侧 14.4mm 的正文列上）。中间的 \quad 在框内是横向空白，
%        显示出来正好是书名与篇名之间的纵向间隔。
%     ③ ★ 内容盒要写成 **与框等宽的定宽盒** `\hbox to \l__zhvt_shukou_chapter_title_height_dim{…\hfil}`。
%        实测：chapter 框是把内容盒沿纵向**居中**摆放的，内容宽度一变，书名就跟着漂：
%        自然宽度的 \hbox（书名 38.4mm）落在 x=157.2mm，加篇名后（57.6mm）落在 x=152.4mm，
%        目录页（无篇名）与正文页的书名就差半格。做成等宽盒后，无论有没有篇名，
%        内容固定从框顶起排，实测两页均落在 x=142.8mm，不再漂。
\ExplSyntaxOn
% 把书口 chapter 框（书名+篇名）起始位置下移 1 字高，避免 4 字卷名顶到鱼尾
% ⚠ cls 的 \__zhvt_calculate_dimensions: 在 \AtBeginDocument 里执行，会重算该 dim；
%   必须把覆写挂到 begin document 之后，否则会被二次覆盖（2026-10-05 实测）。
% (旧方案已移到导言区，直接修改 chapter 框起始位置)
% 卷标题命令 #1=书口篇名(鱼尾) #2=正文卷名（含卷序，卷名内空格用 ~）
%
% ★★ 目录页两项修订（2026-10-05，愚大人指出）★★
%   ① 页码起点统一对齐：
%      原状 —— 目录项写成 \par\hspace{4em}卷名\hspace{4em}页码，页码起点 = 卷名实宽 + 4em。
%      卷名实宽随字数变：「卷一~梁惠王上」6 字 +~、「卷七~离娄上」5 字 +~，
%      差 1 个字位 → 卷七~卷十 4 条的页码整体比其余 10 条高 27.21pt（一个字），看着上下不齐。
%      改法 —— 卷名补全角空格（宽 = 1 字 = 27.215pt）补足到 **7 字符位**再写进目录，
%      14 条页码起点遂全部落在同一条竖线上（起点对齐；末字随页码字数自然参差）。
%      ⚠ 只对目录里的副本 \l_tmpc_tl 补，正文卷名仍用 \l_tmpa_tl 原样 —— 正文像素不受影响。
%   ② 卷名整体上移半字位：前导 \hspace 由 4em 改 3.5em。
%      实测原状目录块在版框内偏下（天头留白 139.6pt / 地脚留白 110.7pt），
%      上移 0.5em = 13.6pt 后上下基本均分，视觉居中。
\NewDocumentCommand { \juan } { mm } {
	\clearpage
	\tl_set:Nn \l__zhvt_chapter_title_tl { #1 }
	\tl_set:Nn \l_tmpa_tl { #2 }
	\tl_set_eq:NN \l_tmpc_tl \l_tmpa_tl
	\int_step_inline:nn { 7 - \str_count:n { #2 } }
	  { \tl_put_right:Nn \l_tmpc_tl { 　 } }
	\tl_set:Nx \l_tmpb_tl { \zhnum { page } }
	\tl_put_left:Nn \l_tmpb_tl { \hspace { 4em } }
	\tl_put_left:Nn \l_tmpb_tl { \l_tmpc_tl }
	\tl_put_left:Nn \l_tmpb_tl { \par \hspace { 3.5em } }
	\addtocontents { toc } { \l_tmpb_tl }
	\skip_horizontal:n { \l__zhvt_font_size_dim * 4 }
	\l_tmpa_tl
	\par
	\thispagestyle { empty }
}
\cs_set_protected:Npn \__zhvt_set_shukou_content: {
	% ★ 鱼尾上方留白（title 静态框）嵌入闲叟钤印：
	%   整页 /Rotate 90，故此处把图用 \rotatebox{90} 转正为竖章；
	%   尺寸卡在版心条幅内（书内约 7mm 宽 × 33.8mm 高，留余量不压鱼尾）。
	\setstaticcontents * { title } {
		\vbox_to_ht:nn { \l__zhvt_baseline_skip_dim } {
			\vskip 14.2pt
			\hbox_to_wd:nn { \l__zhvt_shukou_title_height_dim } {
				\hfil
				\rotatebox [ origin = c ] { 90 } { \includegraphics [ width = 7mm , height = 33.8mm ] { yshs.png } }
				\hfil
			}
		}
	}
% ★ 版心（书口）统一改法（2026-11-16 愚大人定）：18pt＋框内居中＋书口文字同版框色
%   ① 原 27.2pt 时 7 字长篇名（如「廉颇蔺相如列传」）末字溢出 chapter 框、压住页码；
%      改 \fontsize{18pt}{18pt}，前后 \hfil 令字串在框内居中（不随篇名字数漂移）。
%   ② \vbox{\vskip 8.2pt} 定量下移，抵消缩字号后栏宽方向 8.2pt 上偏（像素实测 y 中心 543.3 = 原版）。
%      ⚠ \raisebox / \vspace* 在 flowfram 框内无效；换字号须同步换补偿：16pt→10.1、19pt→7.4。
%   ③ chapter 与 page~number 均置 \color{sealbrown}（HTML 874434），与版框／界行同色；正文仍黑。
%   ⚠ 页码字位、钤印、鱼尾、界行、正文一律未动；五部合订本同此口径。
	\setdynamiccontents * { chapter }
	{ \color{sealbrown} \vbox { \vskip 8.2pt \hbox to \l__zhvt_shukou_chapter_title_height_dim
		{ \hfil \fontsize{18pt}{18pt}\selectfont \l__zhvt_title_str \quad \l__zhvt_chapter_title_tl \hfil } } }
	\setdynamiccontents * { page~number } { \color{sealbrown} \quad \zhnum { page } }
}
\__zhvt_set_shukou_content:
\ExplSyntaxOff

\maketitle

% 单篇，无多卷目录，不渲染目录页

\juan{南越列传}{卷一百十三~南越列传第五十三}
\hspace*{2em}南越王尉佗者。真定人也。姓赵氏。秦时已并天下。略定杨越。置桂林。南海。象郡。以谪徙民。与越杂处十三岁。佗。秦时用为南海龙川令。至二世时。南海尉任嚣病且死。召龙川令赵佗语曰。﹁闻陈胜等作乱。秦为无道。天下苦之。项羽。刘季。陈胜。吴广等州郡各共兴军聚众。虎争天下。中国扰乱。未知所安。豪杰畔秦相立。南海僻远。吾恐盗兵侵地至此。吾欲兴兵绝新道。自备。待诸侯变。会病甚。且番禺负山险。阻南海。东西数千里。颇有中国人相辅。此亦一州之主也。可以立国。郡中长吏无足与言者。故召公告之。﹂即被佗书。行南海尉事。嚣死。佗即移檄告横浦。阳山。湟谿关曰。﹁盗兵且至。急绝道聚兵自守。﹂因稍以法诛秦所置长吏。以其党为假守。秦已破灭。佗即击并桂林。象郡。自立为南越武王。高帝已定天下。为中国劳苦。故释佗弗诛。汉十一年。遣陆贾因立佗为南越王。与剖符通使。和集百越。毋为南边患害。与长沙接境。\\
\hspace*{2em}高后时。有司请禁南越关市铁器。佗曰。﹁高帝立我。通使物。今高后听谗臣。别异蛮夷。隔绝器物。此必长沙王计也。欲倚中国。击灭南越而并王之。自为功也。﹂於是佗乃自尊号为南越武帝。发兵攻长沙边邑。败数县而去焉。高后遣将军隆虑侯灶往击之。会暑湿。士卒大疫。兵不能逾岭。岁馀。高后崩。即罢兵。佗因此以兵威边。财物赂遗闽越。西瓯。骆。役属焉。东西万馀里。乃乘黄屋左纛。称制。与中国侔。\\
\hspace*{2em}及孝文帝元年。初镇抚天下。使告诸侯四夷从代来即位意。喻盛德焉。乃为佗亲冢在真定。置守邑。岁时奉祀。召其从昆弟。尊官厚赐宠之。诏丞相陈平等举可使南越者。平言好畤陆贾。先帝时习使南越。乃召贾以为太中大夫。往使。因让佗自立为帝。曾无一介之使报者。陆贾至南越。王甚恐。为书谢。称曰。﹁蛮夷大长老夫臣佗。前日高后隔异南越。窃疑长沙王谗臣。又遥闻高后尽诛佗宗族。掘烧先人冢。以故自弃。犯长沙边境。且南方卑湿。蛮夷中间。其东闽越千人众号称王。其西瓯骆裸国亦称王。老臣妄窃帝号。聊以自娱。岂敢以闻天王哉。﹂乃顿首谢。原长为籓臣。奉贡职。於是乃下令国中曰。﹁吾闻两雄不俱立。两贤不并世。皇帝。贤天子也。自今以后。去帝制黄屋左纛。﹂陆贾还报。孝文帝大说。遂至孝景时。称臣。使人朝请。然南越其居国窃如故号名。其使天子。称王朝命如诸侯。至建元四年卒。\\
\hspace*{2em}佗孙胡为南越王。此时闽越王郢兴兵击南越边邑。胡使人上书曰。﹁两越俱为籓臣。毋得擅兴兵相攻击。今闽越兴兵侵臣。臣不敢兴兵。唯天子诏之。﹂於是天子多南越义。守职约。为兴师。遣两将军往讨闽越。兵未逾岭。闽越王弟馀善杀郢以降。於是罢兵。\\
\hspace*{2em}天子使庄助往谕意南越王。胡顿首曰。﹁天子乃为臣兴兵讨闽越。死无以报德。﹂遣太子婴齐入宿卫。谓助曰。﹁国新被寇。使者行矣。胡方日夜装入见天子。﹂助去后。其大臣谏胡曰。﹁汉兴兵诛郢。亦行以惊动南越。且先王昔言。事天子期无失礼。要之不可以说好语入见。入见则不得复归。亡国之势也。﹂於是胡称病。竟不入见。后十馀岁。胡实病甚。太子婴齐请归。胡薨。谥为文王。\\
\hspace*{2em}婴齐代立。即藏其先武帝玺。婴齐其入宿卫在长安时。取邯郸樛氏女。生子兴。及即位。上书请立樛氏女为后。兴为嗣。汉数使使者风谕婴齐。婴齐尚乐擅杀生自恣。惧入见要用汉法。比内诸侯。固称病。遂不入见。遣子次公入宿卫。婴齐薨。谥为明王。\\
\hspace*{2em}太子兴代立。其母为太后。太后自未为婴齐姬时。尝与霸陵人安国少季通。及婴齐薨后。元鼎四年。汉使安国少季往谕王。王太后以入朝。比内诸侯。令辩士谏大夫终军等宣其辞。勇士魏臣等辅其缺。卫尉路博德将兵屯桂阳。待使者。王年少。太后中国人也。尝与安国少季通。其使复私焉。国人颇知之。多不附太后。太后恐乱起。亦欲倚汉威。数劝王及群臣求内属。即因使者上书。请比内诸侯。三岁一朝。除边关。於是天子许之。赐其丞相吕嘉银印。及内史。中尉。太傅印。馀得自置。除其故黥劓刑。用汉法。比内诸侯。使者皆留填抚之。王。王太后饬治行装重赍。为入朝具。\\
\hspace*{2em}其相吕嘉年长矣。相三王。宗族官仕为长吏者七十馀人。男尽尚王女。女尽嫁王子兄弟宗室。及苍梧秦王有连。其居国中甚重。越人信之。多为耳目者。得众心愈於王。王之上书。数谏止王。王弗听。有畔心。数称病不见汉使者。使者皆注意嘉。势未能诛。王。王太后亦恐嘉等先事发。乃置酒。介汉使者权。谋诛嘉等。使者皆东乡。太后南乡。王北乡。相嘉。大臣皆西乡。侍坐饮。嘉弟为将。将卒居宫外。酒行。太后谓嘉曰。﹁南越内属。国之利也。而相君苦不便者。何也。﹂以激怒使者。使者狐疑相杖。遂莫敢发。嘉见耳目非是。即起而出。太后怒。欲鏦嘉以矛。王止太后。嘉遂出。分其弟兵就舍。称病。不肯见王及使者。乃阴与大臣作乱。王素无意诛嘉。嘉知之。以故数月不发。太后有淫行。国人不附。欲独诛嘉等。力又不能。\\
\hspace*{2em}天子闻嘉不听王。王。王太后弱孤不能制。使者怯无决。又以为王。王太后已附汉。独吕嘉为乱。不足以兴兵。欲使庄参以二千人往使。参曰。﹁以好往。数人足矣。以武往。二千人无足以为也。﹂辞不可。天子罢参也。郏壮士故济北相韩千秋奋曰。﹁以区区之越。又有王。太后应。独相吕嘉为害。原得勇士二百人。必斩嘉以报。﹂於是天子遣千秋与王太后弟樛乐将二千人往。入越境。吕嘉等乃遂反。下令国中曰。﹁王年少。太后。中国人也。又与使者乱。专欲内属。尽持先王宝器入献天子以自媚。多从人。行至长安。虏卖以为僮仆。取自脱一时之利。无顾赵氏社稷。为万世虑计之意。﹂乃与其弟将卒攻杀王。太后及汉使者。遣人告苍梧秦王及其诸郡县。立明王长男越妻子术阳侯建德为王。而韩千秋兵入。破数小邑。其后越直开道给食。未至番禺四十里。越以兵击千秋等。遂灭之。使人函封汉使者节置塞上。好为谩辞谢罪。发兵守要害处。於是天子曰。﹁韩千秋虽无成功。亦军锋之冠。﹂封其子延年为成安侯。樛乐。其姊为王太后。首原属汉。封其子广德为龙亢侯。乃下赦曰。﹁天子微。诸侯力政。讥臣不讨贼。今吕嘉。建德等反。自立晏如。令罪人及江淮以南楼船十万师往讨之。﹂\\
\hspace*{2em}元鼎五年秋。卫尉路博德为伏波将军。出桂阳。下汇水。主爵都尉杨仆为楼船将军。出豫章。下横浦。故归义越侯二人为戈船。下厉将军。出零陵。或下离水。或柢苍梧。使驰义侯因巴蜀罪人。发夜郎兵。下牂柯江。咸会番禺。\\
\hspace*{2em}元鼎六年冬。楼船将军将精卒先陷寻陕。破石门。得越船粟。因推而前。挫越锋。以数万人待伏波。伏波将军将罪人。道远。会期后。与楼船会乃有千馀人。遂俱进。楼船居前。至番禺。建德。嘉皆城守。楼船自择便处。居东南面。伏波居西北面。会暮。楼船攻败越人。纵火烧城。越素闻伏波名。日暮。不知其兵多少。伏波乃为营。遣使者招降者。赐印。复纵令相招。楼船力攻烧敌。反驱而入伏波营中。犁旦。城中皆降伏波。吕嘉。建德已夜与其属数百人亡入海。以船西去。伏波又因问所得降者贵人。以知吕嘉所之。遣人追之。以其故校尉司马苏弘得建德。封为海常侯。越郎都稽得嘉。封为临蔡侯。\\
\hspace*{2em}苍梧王赵光者。越王同姓。闻汉兵至。及越揭阳令定自定属汉。越桂林监居翁谕瓯骆属汉。皆得为侯。戈船。下厉将军兵及驰义侯所发夜郎兵未下。南越已平矣。遂为九郡。伏波将军益封。楼船将军兵以陷坚为将梁侯。\\
\hspace*{2em}自尉佗初王后。五世九十三岁而国亡焉。\\
\hspace*{2em}太史公曰。尉佗之王。本由任嚣。遭汉初定。列为诸侯。隆虑离湿疫。佗得以益骄。瓯骆相攻。南越动摇。汉兵临境。婴齐入朝。其后亡国。徵自樛女。吕嘉小忠。令佗无后。楼船从欲。怠傲失惑。伏波困穷。智虑愈殖。因祸为福。成败之转。譬若纠墨。\\

% ===== 索隐述赞（司马贞《史记索隐》）· 统一样式 =====
% 题头「索隐述赞曰」五字并入反白框；正文每四句一列，定宽盒 \hfil 居列中；
% \fontsize{24}{41} 列距对齐全书界行（121+41k）；\clearpage 使述赞独占一页。
\clearpage
\shuzanhead{索隐述赞曰}\linebreak
{\shuzanbody\fontsize{24}{41}\selectfont
\hbox to \hsize{\hfil 中原鹿走。群雄莫制。汉事西驱。越权南裔。\hfil}\linebreak
\hbox to \hsize{\hfil 陆贾骋说。尉他去帝。嫪后内朝。吕嘉狼戾。\hfil}\linebreak
\hbox to \hsize{\hfil 君臣不协。卒从剿弃。\hfil}%
\par}
\end{document}
